% ==========================================
% BAB IV DESAIN KONSEP SOLUSI
% ==========================================
\chapter{DESAIN KONSEP SOLUSI}
\label{chap:desain-konsep-solusi}
% Ilustrasikan desain konsep solusi dalam bentuk model konseptual dan penjelasan secara ringkas, 
% beserta perbedaannya dengan sistem saat ini. Ilustrasi harus dapat dibandingkan (\textit{before} and \textit{after}). 
% Karena masih berupa proposal, bab ini hanya berisi gambar desain konsep solusi tersebut dan 
% penjelasan perbandingannya dengan gambar sistem yang ada saat ini (yang tergambar di awal Bab \ref{chap:analisis-masalah}).

% ============================================================================================
% BAB IV DESAIN KONSEP SOLUSI
% ============================================================================================
Bab ini menjabarkan rancangan terperinci dari solusi terpilih pada Bab \ref{chap:analisis-masalah} untuk menjawab rumusan masalah (RM1--RM3) dan menguji hipotesis (H1--H5) pada Bab \ref{chap:pendahuluan}. Pembahasan mencakup arsitektur model konseptual solusi yang diusulkan beserta perbandingannya terhadap alur yang ada, rancangan pemrosesan data dan pembentukan fitur, perancangan model beserta objektif gabungan dan sapuan bobot $\beta$, desain kontrol eksperimen, serta rancangan evaluasi tugas hilir dan analisis representasi internal.

% ============================================================================================
\section{Arsitektur dan Model Konseptual Solusi yang Diusulkan}
\label{sec:ds_arsitektur}

Berdasarkan hasil pemilihan solusi pada Subbab \ref{sub:am_penentuan}, kerangka eksperimen dirancang sebagai pengembangan langsung dari model konseptual alur \textit{brain-tuning} yang ada (Subbab \ref{sub:am_konseptual}). Tata letak lima subsistem utama tetap dipertahankan agar perbandingan antara kondisi sebelum dan sesudah perancangan dapat diamati secara langsung, dengan penambahan satu subsistem baru, yaitu \textbf{Subsistem Analisis Internal}, serta pembaruan mekanisme pada subsistem penyelarasan, kontrol, dan evaluasi.

Gambar \ref{fig:ds_konseptual} mengilustrasikan model konseptual alur \textit{brain-tuning} yang diusulkan. Komponen dan alur yang merupakan pembaruan atau penambahan ditandai dengan garis putus-putus berwarna biru, meliputi: (1) tiga kondisi pelatihan dengan kontrol terstandar, (2) pembaruan berbasis LoRA hingga lapisan target $\ell^*$ disertai kepala proyeksi linear $W_p$ dan \textit{loss} gabungan ($\mathcal{L}_{\text{LM}} + \beta\,\mathcal{L}_{\text{otak}}$), (3) mekanisme sapuan nilai $\beta$ untuk pemetaan keseimbangan, serta (4) subsistem analisis internal yang mencakup perubahan CKA antar-lapisan, skor hierarki kortikal, dan dekomposisi variansi fitur tingkat rendah.

\begin{figure}[h]
  \centering
  \caption{Model konseptual \textit{brain-tuning} yang diusulkan (garis putus-putus berwarna adalah komponen baru dan pembaruan terhadap alur yang ada)}
  \label{fig:ds_konseptual}
  \resizebox{\textwidth}{!}{%
    \begin{tikzpicture}[
        node distance=0.8cm and 0.9cm,
        every node/.style={font=\scriptsize},
        existing/.style={
        rectangle, rounded corners=3pt, draw=black!70, thick, fill=gray!6,
        text width=3.85cm, minimum height=2.35cm, align=left, inner sep=5pt
        },
        proposed/.style={
        rectangle, rounded corners=3pt, draw=blue!75!black, dashed, very thick, fill=blue!5,
        text width=3.85cm, minimum height=2.35cm, align=left, inner sep=5pt
        },
        subbox/.style={
        rectangle, rounded corners=2pt, draw=blue!70!black, dashed, thick, fill=white,
        inner sep=2.5pt, font=\fontsize{6.5}{7.5}\selectfont, align=left
        },
        arr/.style={-{Stealth[length=2mm]}, thick, draw=black!75},
        arrnew/.style={-{Stealth[length=2mm]}, thick, dashed, draw=blue!75!black}
    ]

        % --- BARIS ATAS: ALUR UTAMA ---
        \node[existing] (s1) {
        \textbf{1. Subsistem Data Stimulus \& Otak}\\[2pt]
        $\bullet$ Dataset \textcite{lebel2023natural} (8 subjek)\\
        $\bullet$ Transkrip berpenanda waktu\\
        $\bullet$ Sinyal BOLD fMRI ($\text{TR}=2{,}0\text{ s}$)\\
        \tikz[baseline]\node[subbox]{Pembagian cerita: 22 latih, 4 val, 1 uji};
        };

        \node[existing, right=of s1] (s2) {
        \textbf{2. Subsistem Pembentukan Fitur}\\[2pt]
        $\bullet$ Ekstraksi token terakhir kata $k$\\
        $\bullet$ \textit{Downsampling} Lanczos ke TR\\
        $\bullet$ Penundaan FIR $t-1\dots t-4$ (Pers. \ref{eq:fir})\\
        \tikz[baseline]\node[subbox]{\textit{Hidden state} lapisan target $\ell^*$};
        };

        \node[proposed, right=of s2] (s3) {
        \textbf{3. Subsistem Penyelarasan}\\[2pt]
        $\bullet$ LoRA atensi lapisan $\ell \le \ell^*$ + $W_p$\\
        $\bullet$ $\mathcal{L}_{\text{LM}} + \beta\,\mathcal{L}_{\text{otak}}$ (Kosinus, NT-Xent; ablasi MSE)\\
        $\bullet$ \textbf{Sapuan bobot $\beta$} (skala log)
        };

        % --- BARIS BAWAH: KONTROL, EVALUASI, ANALISIS INTERNAL ---
        \node[proposed, below=1.05cm of s1] (s5) {
        \textbf{5. Subsistem Kontrol}\\[2pt]
        \textbf{Tiga kondisi setara komputasi:}\\
        (1) Bahasa saja ($\beta = 0$)\\
        (2) fMRI asli ($\beta > 0$)\\
        (3) fMRI acak (\textit{TR-} \& \textit{block-shuffled})
        };

        \node[existing, below=1.05cm of s2] (s4) {
        \textbf{4. Subsistem Evaluasi}\\[2pt]
        $\bullet$ \textit{Brain score} ($\rho_v / c_v$, Pers. \ref{eq:pearson})\\
        $\bullet$ Semantik: STS-B, WiC, RTE/MNLI, HellaSwag\\
        $\bullet$ Umum: MMLU, \textit{Perplexity}
        };

        \node[proposed, below=1.05cm of s3] (s6) {
        \textbf{6. Subsistem Analisis Internal}\\[2pt]
        $\bullet$ $\Delta s_{\ell,q}$ per lapisan \& ROI korteks\\
        $\bullet$ Perubahan $1-\text{CKA}$ ($\ell \le \ell^*$ vs $> \ell^*$)\\
        $\bullet$ Skor hierarki $H$ (eksploratori)\\
        $\bullet$ Dekomposisi variansi fitur rendah ($f$)
        };

        % --- PANAH INTERAKSI ---
        \draw[arr] (s1.east) -- node[above, font=\fontsize{6}{7}\selectfont]{Teks} (s2.west);
        \draw[arr] (s2.east) -- node[above, font=\fontsize{6}{7}\selectfont]{$X_t^{\text{delay}}$} (s3.west);

        % Sinyal BOLD dari S1 ke S3 melalui celah tengah
        \draw[arr] ([yshift=-3pt]s1.south east) -- ++(0,-0.32) -| node[pos=0.25, above, font=\fontsize{6}{7}\selectfont]{Sinyal BOLD $y_t$} ([xshift=-12pt]s3.south);

        % S5 (Kontrol) ke S3 (Penyelarasan)
        \draw[arrnew] (s5.north east) -- ++(0,0.42) -| node[pos=0.25, below, font=\fontsize{6}{7}\selectfont, text=blue!75!black]{3 kondisi} ([xshift=-26pt]s3.south);

        % S3 ke S4 (Evaluasi)
        \draw[arr] ([xshift=-8pt]s3.south) -- ++(0,-0.48) -| node[pos=0.25, below, font=\fontsize{6}{7}\selectfont]{Model terlatih} (s4.north);

        % S3 ke S6 (Analisis Internal)
        \draw[arrnew] ([xshift=16pt]s3.south) -- node[right, font=\fontsize{6}{7}\selectfont, text=blue!75!black]{$X_\ell^{\text{tuned}}$} ([xshift=16pt]s6.north);

        % S4 ke S6
        \draw[arrnew] (s4.east) -- node[above, font=\fontsize{6}{7}\selectfont, text=blue!75!black]{$s_{\ell,q}, \hat{W}$} (s6.west);

        % --- LEGENDA ---
        \node[below=0.25cm of s4, font=\fontsize{6.5}{7.5}\selectfont] (leg) {
        \tikz[baseline=-0.5ex]\draw[black!70, thick, fill=gray!6] (0,0) rectangle (0.35,0.2); Komponen eksisting \quad
        \tikz[baseline=-0.5ex]\draw[blue!75!black, dashed, very thick, fill=blue!5] (0,0) rectangle (0.35,0.2); Komponen baru / dimodifikasi
        };

    \end{tikzpicture}
  }
\end{figure}

Secara fungsional, interaksi antarsubsistem bergerak dari penyiapan pasangan stimulus--respons hingga pembuktian hipotesis secara kuantitatif. Pembandingan rinci antara kondisi sebelum dan rancangan perubahan pada tiap subsistem beserta keterkaitannya dengan rumusan masalah (RM) dan hipotesis (H) dirangkum pada Tabel \ref{tab:ds_perbandingan_subsistem}.

\begin{table}[h]
  \centering
  \caption{Perbandingan alur sebelum dan rancangan solusi yang diusulkan per subsistem}
  \label{tab:ds_perbandingan_subsistem}
  \small
  \begin{tabularx}{\textwidth}{p{2.5cm} | p{3.7cm} | p{4.6cm} | X}
    \hline
    \textbf{Subsistem} & \textbf{Kondisi Sebelum (Eksisting)} & \textbf{Perubahan / Rancangan Solusi} & \textbf{RM / H} \\
    \hline
    1. Data stimulus \& sinyal otak & Rekaman fMRI naturalistik dan transkrip berpenanda waktu; pembagian data beragam antarstudi. & Menggunakan dataset \textcite{lebel2023natural} dengan pemisahan ketat berbasis cerita (22 latih, 4 validasi, 1 uji) serta filter \textit{noise ceiling} $c_v > 0{,}1$; \textcite{nastase2021narratives} sebagai validasi opsional. & RM1 \\
    \hline
    2. Pembentukan fitur & Ekstraksi \textit{hidden state} dari seluruh lapisan atau lapisan tetap, lalu \textit{downsampling} dan FIR \textit{delay}. & Penentuan lapisan target $\ell^*$ secara empiris dari kandidat lapisan tengah model awal; pengambilan token terakhir per kata sebelum \textit{downsampling} Lanczos dan FIR 4 TR (Persamaan \ref{eq:fir}). & RM1, H3 \\
    \hline
    3. Penyelarasan (\textit{brain-tuning}) & \textit{Full fine-tuning} pada model kecil atau LoRA pada lapisan tengah dengan satu objektif dan bobot tetap. & Pemasangan LoRA pada proyeksi atensi lapisan $\ell \le \ell^*$ beserta kepala proyeksi linear $W_p$; objektif gabungan (Persamaan \ref{eq:hybrid}) dengan Kosinus dan NT-Xent sebagai utama (MSE ablasi sekunder) serta sapuan bobot $\beta$. & RM1, RM2, H2 \\
    \hline
    4. Evaluasi & Terbatas pada skor \textit{encoding} global atau rata-rata tugas hilir tanpa pemisahan tugas bervalidasi kecil. & Pemisahan evaluasi pemahaman semantik (STS-B, WiC, RTE/MNLI via \textit{probing} 5 \textit{seed}; HellaSwag) dari fungsionalitas umum (MMLU, \textit{Perplexity}); pelaporan skor per tugas. & RM2, H1, H2 \\
    \hline
    5. Kontrol & Sebagian studi memakai kontrol bahasa saja atau pengacakan sinyal secara terpisah. & Penyetaraan tiga kondisi pelatihan (bahasa saja $\beta=0$, fMRI asli $\beta>0$, serta fMRI acak \textit{TR-shuffled} dan \textit{block-shuffled}) dengan teks, langkah, dan \textit{seed} identik. & RM2, H1, H5 \\
    \hline
    6. Analisis internal \textit{(Baru)} & Belum tersedia pada alur \textit{brain-tuning} LLM teks \textit{decoder}. & Kuantifikasi perubahan per lapisan dan ROI ($\Delta s_{\ell,q}$), pergeseran representasi $1-\text{CKA}$ ($\ell \le \ell^*$ vs $> \ell^*$), skor hierarki $H$, dan dekomposisi variansi fitur tingkat rendah $f$. & RM3, H3, H4, H5 \\
    \hline
  \end{tabularx}
\end{table}

% ============================================================================================
\section{Rancangan Data dan Pemrosesan Fitur Stimulus--Otak}
\label{sec:ds_data}

Rancangan data memastikan keselarasan temporal antara stimulus teks dan sinyal BOLD fMRI, dan mencegah bocornya informasi antarpartisi eksperimen (KNF-04).

\begin{enumerate}
  \item \textbf{Partisi data berbasis cerita.} Sebanyak 27 cerita siniar pada dataset \textcite{lebel2023natural} dipisahkan secara utuh pada tingkat cerita menjadi 22 cerita latih, 4 cerita validasi (untuk pemilihan lapisan target $\ell^*$, koefisien regulasi \textit{ridge} $\lambda$, dan kalibrasi rentang $\beta$), serta 1 cerita uji yang direkam berulang pada 5 sesi. Partisi tingkat cerita wajib diterapkan karena penundaan hemodinamik FIR menggabungkan 4 TR berurutan sehingga pemisahan acak per TR akan menimbulkan kebocoran korelasi temporal.
  \item \textbf{Seleksi \textit{voxel} dan pemetaan ROI.} Untuk setiap partisipan ($N=8$), reliabilitas setiap \textit{voxel} dihitung melalui \textit{noise ceiling} $c_v$ sebagai rerata korelasi Pearson berpasangan antar-ulangan pada cerita uji 5 sesi. Hanya \textit{voxel} dengan $c_v > 0{,}1$ yang dipertahankan (dengan analisis sensitivitas pada ambang $0{,}05$ dan $0{,}2$). Himpunan \textit{voxel} ini dikunci sebelum pelatihan dan dikelompokkan ke dalam wilayah minat (\textit{Region of Interest}, ROI) sesuai hierarki pemrosesan bahasa kortikal pada Subbab \ref{sub:sl_otak}: korteks auditori primer sebagai pembanding tingkat rendah, serta area bahasa tingkat tinggi yang meliputi STS, MTG/ITG, ATL, dan IFG.
  \item \textbf{Alur pembentukan fitur stimulus.} Transkrip cerita berpenanda waktu diumpankan ke LLM dalam jendela konteks kausal. Untuk setiap kata $k$, diambil \textit{hidden state} token subkata terakhir $h_k^{(\ell^*)}$ pada lapisan target $\ell^*$ karena token terakhir pada arsitektur \textit{decoder} merangkum konteks kausal sebelumnya. Deret fitur tingkat kata tersebut kemudian diturunkan ke laju akuisisi fMRI ($\text{TR} = 2{,}0\text{ detik}$) menggunakan konvolusi filter Lanczos ($0{,}25\text{ Hz}$), lalu dikonstruksi menjadi vektor fitur berpenundaan $X_t^{\text{delay}} = [x_{t-1};\, x_{t-2};\, x_{t-3};\, x_{t-4}]$ sesuai Persamaan \ref{eq:fir} untuk menangkap jeda hemodinamik 2--8 detik.
\end{enumerate}

% ============================================================================================
\section{Rancangan Model, Lapisan Target, dan Objektif Gabungan}
\label{sec:ds_model_objektif}

\subsection{Konfigurasi Model Dasar dan Adapter PEFT}
\label{sub:ds_model}

Eksperimen difokuskan pada LLM teks \textit{decoder-only} berukuran menengah dengan model utama \textbf{Qwen3-4B} (4B parameter) dan model pembanding \textbf{Llama-3-8B} atau \textbf{Llama-3.1-8B} (8B parameter), serta \textbf{Mistral-7B} sebagai model opsional (KNF-08). Seluruh model dilatih menggunakan mekanisme PEFT (Subbab \ref{sub:sl_finetuning}) agar tetap berjalan dalam batas memori VRAM 16 GB (KNF-01):
\begin{itemize}
  \item \textbf{LoRA (presisi 16-bit):} diterapkan sebagai konfigurasi utama pada model berukuran hingga 4B parameter, dengan modul adapter pada matriks proyeksi modul \textit{self-attention} ($W^Q, W^K, W^V, W^O$) sesuai Persamaan \ref{eq:multihead} dan \ref{eq:lora}.
  \item \textbf{QLoRA (kuantisasi 4-bit NormalFloat):} disiapkan sebagai opsi pembaruan untuk model berukuran 4--8B parameter \autocite{dettmers2023qloraefficientfinetuningquantized} jika kebutuhan memori dan gradien dalam presisi 16-bit melampaui kapasitas 16 GB.
\end{itemize}

\subsection{Pemilihan Lapisan Target ($\ell^*$) dan Kepala Proyeksi}
\label{sub:ds_lapisan}

Penentuan lapisan target dilakukan dalam dua tahap sebelum pelatihan utama:
\begin{enumerate}
  \item \textbf{Pemindaian lapisan pada model awal.} Model awal (\textit{pretrained} tanpa \textit{tuning}) dievaluasi pada cerita validasi menggunakan \textit{voxelwise encoding model} (Persamaan \ref{eq:ridge}--\ref{eq:pearson}) untuk tiap lapisan $\ell \in \{1, \dots, L\}$. Lapisan $\ell^*$ dipilih dari kandidat sepertiga hingga dua pertiga kedalaman jaringan (lapisan tengah) yang memperoleh rerata \textit{brain score} validasi tertinggi, selaras dengan hierarki representasi semantik pada \textcite{mischler2024contextual} dan \textcite{xiao2026beyond}.
  \item \textbf{Cakupan pemasangan LoRA ($\ell \le \ell^*$).} Selama \textit{brain-tuning}, fitur berpenundaan $X_t^{\text{delay}}$ diekstrak dari lapisan $\ell^*$ dan dipetakan ke ruang \textit{voxel} partisipan melalui kepala proyeksi linear $\hat{y}_t = W_p X_t^{\text{delay}} + b$ (Persamaan \ref{eq:proyeksi}). Karena gradien dari \textit{loss} otak hanya merambat mundur dari lapisan $\ell^*$ ke lapisan-lapisan di bawahnya, adapter LoRA dipasang pada proyeksi atensi seluruh lapisan dari lapisan pertama hingga lapisan $\ell^*$ ($\ell \le \ell^*$). Lapisan di atas $\ell^*$ ($\ell > \ell^*$) dibiarkan tanpa adapter sehingga perubahan representasinya murni mencerminkan propagasi tak langsung, yang kemudian diuji melalui H3. Kepala proyeksi $W_p$ dilatih bersama parameter LoRA selama fase \textit{brain-tuning} dan dilepas saat fase evaluasi hilir maupun pengujian \textit{encoding} independen.
\end{enumerate}

\subsection{Objektif Gabungan dan Sapuan Bobot $\beta$}
\label{sub:ds_objektif}

Pembaruan parameter adapter LoRA dan kepala proyeksi $W_p$ dikendalikan oleh fungsi objektif gabungan pada Persamaan \ref{eq:hybrid}:
\begin{equation*}
  \mathcal{L}_{\text{total}} = \mathcal{L}_{\text{LM}} + \beta\, \mathcal{L}_{\text{otak}}
\end{equation*}
dengan $\mathcal{L}_{\text{LM}}$ berupa \textit{cross-entropy loss} prediksi token berikutnya pada transkrip cerita latih untuk mempertahankan kompetensi bahasa model. Untuk komponen $\mathcal{L}_{\text{otak}}$, digunakan dua fungsi objektif utama yang invarian terhadap skala respons \textit{voxel}:
\begin{enumerate}
  \item \textbf{Cosine Loss ($\mathcal{L}_{\text{cos}}$, Persamaan \ref{eq:cosine}):} memaksimalkan keselarasan arah vektor antara prediksi $\hat{y}_i$ dan respons BOLD target $y_i$ pada setiap langkah waktu dengan rentang nilai berbatas $[0, 2]$.
  \item \textbf{NT-Xent Loss ($\mathcal{L}_{\text{NT-Xent}}$, Persamaan \ref{eq:ntxent}):} memberikan pembelajaran kontrastif dengan memperlakukan pasangan $(\hat{y}_i, y_i)$ pada TR yang bersesuaian sebagai pasangan positif dan pasangan pada TR lain dalam \textit{mini-batch} sebagai sampel negatif, dengan parameter temperatur $\tau$ (nilai awal $\tau = 0{,}1$ yang dikalibrasi pada uji pendahuluan).
\end{enumerate}
Sementara itu, \textbf{Mean Squared Error ($\mathcal{L}_{\text{MSE}}$, Persamaan \ref{eq:mse})} ditempatkan sebagai objektif sekunder (ablasi pembanding) apabila anggaran komputasi memungkinkan setelah seluruh eksperimen utama selesai.

Untuk memetakan keseimbangan antara keselarasan otak, peningkatan pemahaman semantik, dan stabilitas fungsionalitas umum model (H2), dilakukan \textbf{sapuan (\textit{sweep}) koefisien $\beta$} pada skala logaritmik yang mencakup titik nol hingga dominasi sinyal otak, yaitu $\beta \in \{0;\, 0{,}01;\, 0{,}05;\, 0{,}1;\, 0{,}5;\, 1{,}0;\, 5{,}0;\, 10{,}0\}$ (rentang akhir disesuaikan dengan hasil uji pendahuluan pada Qwen3-4B). Pada setiap titik $\beta$, dicatat \textit{brain score} uji, rerata skor tugas semantik, akurasi MMLU, dan \textit{perplexity} untuk mengidentifikasi rentang nilai $\beta$ optimal sebelum penurunan fungsionalitas terjadi.

% ============================================================================================
\section{Desain Kondisi Kontrol Eksperimen}
\label{sec:ds_kontrol}

Dirancang 3 kondisi pelatihan dengan komputasi, korpus transkrip, \textit{rank} LoRA, laju belajar, ukuran \textit{batch}, langkah pembaruan, dan \textit{random seed} setara (RM2, KNF-03):

\begin{enumerate}
  \item \textbf{Kondisi Kontrol Bahasa Saja (\textit{Language-Only Fine-Tuning}, $\beta = 0$).} Model hanya dilatih menggunakan $\mathcal{L}_{\text{LM}}$ pada transkrip 22 cerita latih tanpa sinyal otak. Kondisi ini berfungsi sebagai tolok ukur untuk mengisolasi pengaruh adaptasi teks naratif terhadap tugas hilir dan pergeseran representasi CKA.
  \item \textbf{Kondisi Perlakuan Utama (\textit{Real fMRI Brain-Tuning}, $\beta > 0$).} Model dilatih dengan objektif gabungan $\mathcal{L}_{\text{LM}} + \beta\,\mathcal{L}_{\text{otak}}$ menggunakan pasangan stimulus dan respons BOLD asli $y_t^{\text{asli}}$ per partisipan.
  \item \textbf{Kondisi Kontrol fMRI Acak (\textit{Shuffled fMRI Brain-Tuning}, $\beta > 0$).} Model dilatih dengan persamaan dan nilai $\beta$ yang persis sama dengan kondisi kedua, tetapi urutan temporal respons BOLD latih diacak sehingga korespondensi antara kata dan aktivitas otak hilang:
  \begin{itemize}
    \item \textit{Pengacakan per TR (\textit{TR-shuffled}, kontrol acak utama):} indeks waktu $t$ pada matriks respons $Y_{\text{latih}}$ dipermutasi secara acak dengan permutasi yang seragam di seluruh \textit{voxel} \autocite{negi2025brain}. Metode ini mempertahankan distribusi marginal nilai BOLD tiap \textit{voxel}, tetapi memutus autokorelasi temporal dan keterkaitan dengan stimulus.
    \item \textit{Pengacakan per blok (\textit{block-shuffled}, kontrol acak sekunder):} respons latih dipotong menjadi segmen berdurasi tetap yang lebih panjang dari jendela penundaan hemodinamik ($> 8\text{ detik}$, misalnya 10 TR atau $20\text{ detik}$) lalu urutan antarblok diacak. Kontrol sekunder ini dijalankan pada model utama dan satu nilai $\beta$ terpilih untuk memverifikasi bahwa simpulan tidak bergantung pada hilangnya struktur autokorelasi lokal akibat pengacakan per TR.
  \end{itemize}
\end{enumerate}

% Seluruh kondisi dievaluasi pada cerita uji menggunakan sinyal fMRI asli yang tidak diacak.

% ============================================================================================
\section{Rancangan Evaluasi Hilir dan Analisis Representasi Internal}
\label{sec:ds_evaluasi_analisis}

Tahap evaluasi dan analisis dirancang secara terstruktur untuk menguji kelima hipotesis penelitian secara kuantitatif.

\subsection{Evaluasi Tugas Hilir dan Fungsionalitas Model (H1 dan H2)}
\label{sub:ds_eval_hilir}

Setelah pelatihan selesai, adapter LoRA digabungkan atau dipertahankan dalam mode inferensi beku tanpa pembaruan bobot lebih lanjut, lalu dievaluasi pada dua kelompok tolok ukur (Tabel \ref{tab:sl_downstream}):
\begin{itemize}
  \item \textbf{Pengujian H1 (Pemahaman Semantik Hilir):} Pada tugas klasifikasi dan regresi pasangan kalimat (STS-B, WiC, serta RTE atau MNLI), digunakan protokol \textit{linear probing} dengan melatih kepala klasifikasi/regresi linear di atas representasi lapisan terakhir pada 5 \textit{random seed} berbeda untuk menekan variansi inisialisasi kepala. Untuk HellaSwag, evaluasi dilakukan secara deterministik menggunakan rerata \textit{log-likelihood} kandidat jawaban (Persamaan \ref{eq:ll}). Skor kondisi fMRI asli dibandingkan secara berpasangan terhadap kondisi $\beta = 0$ dan kondisi fMRI acak.
  \item \textbf{Pengujian H2 (Pemetaan Keseimbangan dan Fungsionalitas Umum):} Pada setiap nilai sapuan $\beta$, kemampuan pengetahuan umum diukur dengan MMLU (Persamaan \ref{eq:ll}) dan kelancaran pemodelan bahasa diukur dengan \textit{Perplexity} pada korpus uji (Persamaan \ref{eq:ppl_w}). Hasilnya diplot bersama kurva kenaikan \textit{brain score} dan skor semantik terhadap $\beta$ untuk membuktikan keberadaan jendela operasi aman.
\end{itemize}

\subsection{Analisis Perubahan Representasi Internal (H3, H4, dan H5)}
\label{sub:ds_analisis_internal}

Untuk menjawab RM3 mengenai mekanisme internal di balik perubahan perilaku model, diterapkan empat analisis kuantitatif pada cerita uji:
\begin{enumerate}
  \item \textbf{Pergeseran Representasi Antar-lapisan via CKA (Pengujian H3).} Untuk setiap lapisan $\ell \in \{1, \dots, L\}$, matriks representasi cerita uji dari model awal ($X_\ell^{\text{awal}}$) dan model pascapelatihan ($X_\ell^{\text{tuned}}$) dibandingkan menggunakan CKA linear (Persamaan \ref{eq:cka}) untuk memperoleh besaran perubahan $\Delta_\ell = 1 - \mathrm{CKA}(X_\ell^{\text{awal}}, X_\ell^{\text{tuned}})$ (Persamaan \ref{eq:delta_cka}). Rerata $\Delta_\ell$ pada kelompok lapisan yang mendapat gradien otak langsung ($\ell \le \ell^*$) dibandingkan terhadap kelompok lapisan di atasnya ($\ell > \ell^*$), serta dibandingkan terhadap nilai $\Delta_\ell$ pada kontrol bahasa saja ($\beta = 0$).
  \item \textbf{Perubahan Keselarasan per ROI dan Skor Hierarki (Pengujian H4).} \textit{Encoding model} independen dilatih ulang pada 22 cerita latih dan diuji pada cerita uji untuk setiap lapisan $\ell$ dan kondisi model. Selisih keselarasan $\Delta s_{\ell^*, q} = s^{\text{tuned}}_{\ell^*, q} - s^{\text{awal}}_{\ell^*, q}$ dihitung per partisipan dan diringkas dengan median per ROI. Kenaikan pada jaringan bahasa tingkat tinggi (STS, MTG/ITG, ATL, IFG) diuji secara berpasangan terhadap kenaikan pada korteks auditori primer. Sebagai analisis eksploratori tambahan, dihitung skor keselarasan hierarki $H = \mathrm{corr}(\mathrm{COM}_q, d_q)$ melalui Persamaan \ref{eq:norm}--\ref{eq:hierarki} untuk mengamati apakah pusat massa lapisan bergeser sejalan dengan jarak kortikal $d_q$.
  \item \textbf{Dekomposisi Variansi Fitur Tingkat Rendah (Pengujian H5).} Menggunakan \textit{banded ridge regression} \autocite{nunez2019banded} dengan penalti terpisah $\lambda_g$ per kelompok fitur, dibangun dua model prediksi BOLD pada cerita uji: model hanya dengan fitur teks tingkat rendah (panjang kata, frekuensi kata, posisi kata dalam kalimat, serta \textit{embedding} statis tanpa konteks) dan model gabungan yang menyertakan \textit{hidden state} lapisan $\ell^*$. Porsi variansi yang dijelaskan fitur tingkat rendah dihitung sebagai $f = R^2_{\text{rendah}} / R^2_{\text{gabungan}}$ (Persamaan \ref{eq:porsi}). Penurunan nilai $f$ setelah \textit{brain-tuning} fMRI asli dibandingkan terhadap perubahan nilai $f$ pada kondisi fMRI acak.
\end{enumerate}

Keterkaitan antara hipotesis, metrik pengukuran, kondisi pembanding, dan kriteria pembuktiannya dirangkum pada Tabel \ref{tab:ds_matriks_hipotesis}. Rancangan konseptual beserta matriks evaluasi pada bab ini menjadi acuan penyusunan lingkungan implementasi, protokol verifikasi dan validasi statistik, jadwal kerja, serta mitigasi risiko pada Bab \ref{chap:rencana-selanjutnya}.

\begin{table}[h]
  \centering
  \caption{Matriks rancangan evaluasi dan pembuktian hipotesis penelitian}
  \label{tab:ds_matriks_hipotesis}
  \small
  \begin{tabularx}{\textwidth}{p{1.2cm} | p{3.6cm} | p{3.3cm} | X}
    \hline
    \textbf{Hipo.} & \textbf{Metrik \& Persamaan} & \textbf{Pembanding} & \textbf{Kriteria Dukungan Hipotesis} \\
    \hline
    H1 & Skor STS-B, WiC, RTE/MNLI (\textit{probing} 5 \textit{seed}), HellaSwag (Persamaan \ref{eq:ll}) & fMRI asli vs bahasa saja ($\beta=0$) dan fMRI acak & Skor tugas semantik pada fMRI asli lebih tinggi secara signifikan ($p < 0{,}05$ terkoreksi FDR) daripada kedua kontrol. \\
    \hline
    H2 & \textit{Brain score} (Persamaan \ref{eq:pearson}), MMLU (Persamaan \ref{eq:ll}), \textit{Perplexity} (Persamaan \ref{eq:ppl_w}) & Kurva metrik lintas sapuan bobot $\beta$ terhadap $\beta=0$ & Ditemukan rentang $\beta > 0$ dengan kenaikan \textit{brain score} tanpa penurunan berarti MMLU/\textit{Perplexity}, sedangkan $\beta$ besar menurunkannya. \\
    \hline
    H3 & Perubahan representasi $\Delta_\ell = 1 - \mathrm{CKA}$ (Persamaan \ref{eq:cka}--\ref{eq:delta_cka}) & Kelompok $\ell \le \ell^*$ vs $\ell > \ell^*$; fMRI asli vs $\beta=0$ & Rerata $\Delta_\ell$ pada $\ell \le \ell^*$ lebih besar secara signifikan daripada $\ell > \ell^*$ dan lebih besar daripada kontrol $\beta=0$. \\
    \hline
    H4 & Selisih \textit{brain score} $\Delta s_{\ell^*, q}$ per ROI; skor hierarki $H$ (Persamaan \ref{eq:norm}--\ref{eq:hierarki}) & ROI tingkat tinggi (STS, MTG/ITG, ATL, IFG) vs auditori primer & $\Delta s_{\ell^*, q}$ pada ROI tingkat tinggi lebih besar secara signifikan daripada korteks auditori primer lintas partisipan. \\
    \hline
    H5 & Porsi variansi fitur rendah $f = R^2_{\text{rendah}} / R^2_{\text{gabungan}}$ (Persamaan \ref{eq:porsi}) & Sebelum vs sesudah \textit{tuning}; fMRI asli vs fMRI acak & Nilai $f$ menurun setelah \textit{brain-tuning} fMRI asli dan penurunannya lebih besar dibandingkan kondisi fMRI acak. \\
    \hline
  \end{tabularx}
\end{table}